% Snapshot of the "Mathematical background" section of main.tex (lines 335-805),
% taken on 2026-11-10 before rewriting that section from scratch.
% Standalone: compile from this folder with `latexmk -pdf math-background_2026-11-10.tex`.
% References to other sections of the survey will show as "??".
\documentclass[11pt]{article}
\usepackage[margin=2.5cm]{geometry}
\usepackage{amsmath,amssymb,amsthm}
\usepackage{enumitem}
\usepackage{xcolor}
\usepackage[numbers,sort&compress]{natbib}
\usepackage{hyperref}
\usepackage{cleveref}

% Theorem-like environments provided by the MDPI class in main.tex
\theoremstyle{plain}
\newtheorem{Theorem}{Theorem}[section]
\newtheorem{Proposition}{Proposition}[section]
\newtheorem{Lemma}{Lemma}[section]
\theoremstyle{definition}
\newtheorem{Definition}{Definition}[section]
\newtheorem{Example}{Example}[section]
\crefname{Theorem}{Theorem}{Theorems}
\crefname{Proposition}{Proposition}{Propositions}
\crefname{Lemma}{Lemma}{Lemmas}
\crefname{Definition}{Definition}{Definitions}
\crefname{Example}{Example}{Examples}

\newcommand{\pnote}[1]{\textcolor{red}{PL: #1}}

\begin{document}

\section{Background: Mathematical Foundations of Symmetry (old version)}\label{sec:background}

% ===== BEGIN copy of main.tex, lines 335-805 =====
% Groups, products, actions, orbits, stabilizers, quotients, symmetry breaking
\subsection{Groups and Group Actions}\label{sec:group-theory-backbone}

The precise language of symmetry is group theory.
A group $G$ models symmetry in the abstract (internal structure), while a group action is the realization of those transformations on a specific space $\mathcal X$.
Actions are necessary because the abstract group $G$ does not dictate a geometry: it can act on a discrete set, a vector space, a smooth manifold, or a function space.

\subsubsection{Groups}\label{subsec:groups}

\begin{Definition}[Group]\label{def:group}
A \emph{group} is a set $G$ with a binary operation $(g,h)\mapsto gh$ satisfying:
\begin{enumerate}[label=(G\arabic*), leftmargin=*, labelsep=0.5em]
    \item \textbf{Closure.} $gh \in G$ for all $g,h\in G$;
    \item \textbf{Associativity.} $(gh)k = g(hk)$ for all $g,h,k\in G$;
    \item \textbf{Identity.} there is $e\in G$ with $eg = ge = g$ for all $g\in G$;
    \item \textbf{Inverses.} every $g\in G$ has a $g^{-1}\in G$ with $gg^{-1} = g^{-1}g = e$.
\end{enumerate}
$G$ is \emph{abelian} if moreover $gh = hg$ for all $g,h\in G$.
\end{Definition}

Closure can be implicit in the binary operation taking values in $G$, but it is the axiom to check first for instance when $G$ is derived from a larger set; relaxing these axioms leads to the weaker structures of \Cref{sec:extent-restriction}.
A \emph{subgroup} $H\leq G$ is a subset containing $e$ and closed under products and inverses; it is \emph{normal}, $H\trianglelefteq G$, if $gHg^{-1} = H$ for all $g\in G$.
A map $\varphi: G\to G'$ is a \emph{group homomorphism} if $\varphi(gh) = \varphi(g)\varphi(h)$, an \emph{isomorphism} ($G\cong G'$) if it is moreover bijective---called \emph{automorphism} when $G=G'$.

\begin{Example}[Groups used in this survey]\label{ex:groups}
\phantom{u}
\begin{itemize}
    \item The \emph{symmetric group} $S_n$ of permutations of $\{1,\dots,n\}$, non-abelian for $n\geq 3$; it can act on sets, on graph nodes and on the hidden units of a layer.
    \item The \emph{cyclic group} $C_n\cong\mathbb{Z}/n\mathbb{Z}$ of rotations by multiples of $2\pi/n$, and the \emph{dihedral group} $D_n$ (order $2n$) which adds the $n$ reflections of a regular $n$-gon; $C_4$ and $D_4$ are the symmetries of a square pixel grid.
    \item The \emph{general linear group} $GL(n,\mathbb{R})$ of invertible $n\times n$ matrices, the \emph{orthogonal group} $O(n) = \{A : A^\top A = I\}$ and its normal subgroup of rotations $SO(n) = \{A\in O(n) : \det A = 1\}$.
    \item The multiplicative group $\mathbb{R}_{>0}$, acting by positive rescaling on the weights of ReLU units (\Cref{sec:parameter-space-symmetries}).
\end{itemize}
\end{Example}

\noindent New groups are built from old ones by products.

\begin{Definition}[Direct and semidirect product]\label{def:semidirect-product}
Let $N,H$ be groups.
The \emph{direct product} $N\times H$ is the set $N\times H$ with componentwise operation $(n_1,h_1)(n_2,h_2) = (n_1n_2,\,h_1h_2)$.
Given an action $\varphi: H\to\operatorname{Aut}(N)$, the \emph{semidirect product} $N\rtimes_\varphi H$ is the set $N\times H$ with
\[
    (n_1,h_1)(n_2,h_2) \;=\; \bigl(n_1\,\varphi(h_1)(n_2),\; h_1h_2\bigr).
\]
Then $N\times\{e\}\trianglelefteq N\rtimes_\varphi H$, $\{e\}\times H$ is a subgroup, and the direct product is the case of trivial $\varphi$.
\end{Definition}

A direct product models two symmetries acting independently; a semidirect product models one symmetry $H$ acting on the other before they compose.

\begin{Example}[Semidirect products]\label{ex:semidirect-products}
\phantom{u}
\begin{itemize}
    \item \textbf{Euclidean group.} $E(n) = \mathbb{R}^n\rtimes O(n)$, with $(v,A)$ acting on $\mathbb{R}^n$ by $x\mapsto Ax+v$ and $\varphi(A)v = Av$.
    Composition follows $(v_1,A_1)(v_2,A_2) = (v_1 + A_1v_2,\; A_1A_2)$ i.e.\ the rotation $A_1$ twists the translation $v_2$ leading to the composed action $x\mapsto A_1(A_2x+v_2)+v_1$.
    The rigid motions form $SE(n) = \mathbb{R}^n\rtimes SO(n)$.
    \item \textbf{Hidden-layer symmetries.} Let a hidden layer of $n$ ReLU units have incoming weights $W_1\in\mathbb{R}^{n\times d}$ and outgoing weights $W_2\in\mathbb{R}^{m\times n}$.
    For a permutation matrix $P_\sigma$ and $D_\lambda = \operatorname{diag}(\lambda_1,\dots,\lambda_n)$ with $\lambda_i>0$, the map $(W_1,W_2)\mapsto(D_\lambda P_\sigma W_1,\; W_2P_\sigma^{\top}D_\lambda^{-1})$ leaves the network function unchanged.
    The matrices $D_\lambda P_\sigma$ form the group $\mathbb{R}_{>0}^n\rtimes S_n$, with $\varphi(\sigma)(\lambda)=P_\sigma \lambda P_\sigma^{\top}$ i.e.\ $\varphi(\sigma)(\lambda)_i = \lambda_{\sigma(i)}$.
\end{itemize}
\end{Example}

\subsubsection{Group Actions}
\label{subsec:group-actions}

\begin{Definition}[Group action]\label{def:group-action}
A (left) \emph{action} of a group $G$ on a set $\mathcal{X}$ is a map $G\times\mathcal{X}\to\mathcal{X}$, $(g,x)\mapsto g\cdot x$, such that
\begin{enumerate}[label=(A\arabic*), leftmargin=*, labelsep=0.5em]
    \item $e\cdot x = x$ for all $x\in\mathcal{X}$;
    \item $g\cdot(h\cdot x) = (gh)\cdot x$ for all $g,h\in G$, $x\in\mathcal{X}$.
\end{enumerate}
Equivalently, an action is a homomorphism $G\to\operatorname{Sym}(\mathcal{X})$, $g\mapsto(x\mapsto g\cdot x)$.
\end{Definition}

\begin{Definition}[Orbit and stabilizer]\label{def:orbit-stabilizer}
For $x\in\mathcal{X}$, the \emph{orbit} of $x$ and the \emph{stabilizer} of $x$ are
\[
    G\cdot x = \{g\cdot x : g\in G\}\subseteq\mathcal{X},
    \qquad
    G_x = \operatorname{Stab}(x) = \{g\in G : g\cdot x = x\}\leq G.
\]
\end{Definition}

The relation $x\sim y \iff y\in G\cdot x$ is an equivalence relation, so orbits partition $\mathcal{X}$.
The \emph{quotient} $\mathcal{X}/G$ is the set of orbits, with projection $\pi: \mathcal{X}\to\mathcal{X}/G$, $x\mapsto G\cdot x$.
It is the space of objects \emph{up to symmetry}, and it links invariance to quotients: by \Cref{def:invariance}, $\phi$ is $G$-invariant iff it is constant on orbits, i.e.\ iff $\phi = \bar\phi\circ\pi$ for a unique $\bar\phi : \mathcal{X}/G\to\mathcal{Y}$.
Parameter-space quotients (\Cref{sec:quotient-repr-align}) and orbit or quotient losses (\Cref{sec:criterion-symmetries}) are instances of this construction.

For $H\leq G$, the \emph{left cosets} $gH = \{gh : h\in H\}$ are the orbits of $H$ acting on $G$ by right multiplication.
They therefore partition $G$, the set of cosets is written $G/H$, and each coset is in bijection with $H$ via $h\mapsto gh$.
For finite $G$ this gives Lagrange's theorem, $|G| = |G/H|\,|H|$.
If $H\trianglelefteq G$, then $G/H$ is itself a group, the \emph{quotient group}, with $(g_1H)(g_2H) = g_1g_2H$ (e.g.\ $O(n)/SO(n)\cong C_2$, recording orientation).

\begin{Theorem}[Orbit--Stabilizer]\label{thm:orbit-stabilizer}
Let $G$ act on $\mathcal{X}$ and $x\in\mathcal{X}$. The map $gG_x\mapsto g\cdot x$ is a well-defined bijection $G/G_x\to G\cdot x$. In particular, if $G$ is finite, $|G\cdot x| = |G|/|G_x|$.
\end{Theorem}
\begin{proof}[Proof sketch]
$g\cdot x = g'\cdot x \iff g^{-1}g'\in G_x \iff gG_x = g'G_x$; this gives well-definedness and injectivity at once, and surjectivity holds by definition of the orbit.
\end{proof}

\begin{Example}[Pixels under rotation]\label{ex:orbit-stabilizer-pixels}
Let $C_4$ act on the pixel positions of an $n\times n$ image, $n$ odd, by rotation about the centre.
The centre pixel has stabilizer $C_4$ and orbit of size $1$; every other pixel has trivial stabilizer and orbit of size $4 = |C_4|/1$.
\end{Example}

\begin{Definition}[Free, transitive and faithful actions]\label{def:free-transitive}
An action of $G$ on $\mathcal{X}$ is
\begin{itemize}
    \item \emph{free} if $G_x = \{e\}$ for all $x$, i.e.\ $g\cdot x = x\implies g=e$;
    \item \emph{transitive} if $\mathcal{X}$ is a single orbit: for all $x,y$ there is $g$ with $g\cdot x = y$;
    \item \emph{faithful} if $G\to\operatorname{Sym}(\mathcal{X})$ is injective, i.e.\ $\bigcap_{x}G_x = \{e\}$.
\end{itemize}
A free and transitive action is \emph{regular} or \emph{simply transitive}, and $\mathcal{X}$ is then a \emph{$G$-torsor}: any basepoint $x_0$ gives a bijection $G\to\mathcal{X}$, $g\mapsto g\cdot x_0$, but no point of $\mathcal{X}$ is distinguished.
\end{Definition}

Free implies faithful; freeness and transitivity are independent.
$S_n$ acts transitively on $\{1,\dots,n\}$ but not freely ($G_1\cong S_{n-1}$).
Conversely, $\mathbb{R}_{>0}$ acts freely but not transitively on the weights $(w_1,w_2)\in(\mathbb{R}^d\setminus\{0\})\times\mathbb{R}^m$ of a single ReLU unit by $\lambda\cdot(w_1,w_2) = (\lambda w_1,\lambda^{-1}w_2)$.
Each orbit is a copy of $\mathbb{R}_{>0}$, a curve along which the network function is constant, and there are infinitely many orbits.

\subsubsection{Stabilizers and symmetry breaking}
\label{subsec:symmetry-breaking}

Stabilizers measure how much symmetry an individual object retains.

\begin{Definition}[Symmetry breaking]\label{def:symmetry-breaking}
Let $G$ act on a set $\mathcal{Y}$ of admissible states, models or ground truths of a problem.
\emph{Symmetry breaking} occurs when the object of interest $y_0\in\mathcal{Y}$ has a stabilizer $G_{y_0}\subsetneq G$ strictly smaller than the symmetry group of the problem.
The distinct but equally valid instances $g\cdot y_0$ then form the orbit $G\cdot y_0\cong G/G_{y_0}$ (\Cref{thm:orbit-stabilizer}).
\end{Definition}

For example, a Hamiltonian invariant under the translations $\mathbb{R}^3$ can have a crystalline ground state $y_0$ invariant only under a lattice $\mathbb{Z}^3$; the equally valid ground states are then indexed by $\mathbb{R}^3/\mathbb{Z}^3$, a 3-torus.
Equivariant maps can only preserve or enlarge stabilizers ($G_x\leq G_{\phi(x)}$, Curie's principle, \Cref{sec:curie-principle}); this is the root of the functional, distributional and parameter-space forms of symmetry breaking discussed in \Cref{sec:symmetry-breaking-dl}.


%
\subsection{Representations}
\label{subsec:representation-theory}

Equivariant architectures act on vector spaces of features, on which groups act \emph{linearly}.

\begin{Definition}[Linear representation]\label{def:linear-representation}
A \emph{representation} of $G$ on a vector space $V$ (over $\mathbb{R}$ or $\mathbb{C}$) is a homomorphism $\rho : G\to GL(V)$, i.e.\ an action of $G$ on $V$ by linear maps, $\rho(g_1g_2) = \rho(g_1)\rho(g_2)$.
It is \emph{unitary} (orthogonal over $\mathbb{R}$) if $V$ carries an inner product preserved by every $\rho(g)$.
A linear map $T : V\to V'$ between representations $\rho,\rho'$ is an \emph{intertwiner}, or \emph{$G$-equivariant linear map}, if
\[
    T\rho(g) = \rho'(g)T \qquad\text{for all } g\in G,
\]
and $\rho\cong\rho'$ are \emph{equivalent} if some intertwiner is invertible.
We write $\operatorname{Hom}_G(V,V')$ for the space of intertwiners.
\end{Definition}

Intertwiners are exactly the equivariant linear layers of \Cref{def:equivariance} with linear actions $\rho_X = \rho$, $\rho_Y = \rho'$.

\begin{Example}[Representations]\label{ex:representations-recap}
\phantom{u}
\begin{itemize}
    \item The \emph{trivial} representation $\rho(g) = \operatorname{id}_V$: intertwiners into it are invariant linear maps.
    \item The \emph{defining} representation of $O(n)$ on $\mathbb{R}^n$, $\rho(A)=A$ (e.g.\ forces or velocities under rotation).
    \item The \emph{permutation} representation of $S_n$ on $\mathbb{R}^n$, $\rho(\sigma)_{ij} = \mathbf{1}[i = \sigma(j)]$. Its intertwiners $\mathbb{R}^n\to\mathbb{R}^n$ are exactly $T = aI + b\mathbf{1}\mathbf{1}^\top$, the linear layer of DeepSets (\Cref{sec:gdl}).
    \item The \emph{regular} representation of a finite group on $\mathbb{R}^{G} = \{f : G\to\mathbb{R}\}$, $(\rho(g)f)(h) = f(g^{-1}h)$: $G$ acting on functions on itself by left translation.
\end{itemize}
\end{Example}

New representations are obtained from $\rho_1: G\to GL(V_1)$ and $\rho_2 : G\to GL(V_2)$ by the \emph{direct sum} $(\rho_1\oplus\rho_2)(g) = \rho_1(g)\oplus\rho_2(g)$ on $V_1\oplus V_2$ (block-diagonal matrices), and the \emph{tensor product} $(\rho_1\otimes\rho_2)(g) = \rho_1(g)\otimes\rho_2(g)$ on $V_1\otimes V_2$ (Kronecker products).
For instance, a rank-$2$ tensor such as a stress or inertia tensor transforms under $O(3)$ by $\rho\otimes\rho$, $M\mapsto AMA^\top$.
Direct sums build feature spaces carrying several types of features; tensor products combine features multiplicatively.

\begin{Definition}[Irreducibility]\label{def:irrep}
A subspace $W\subseteq V$ is \emph{invariant} if $\rho(g)W\subseteq W$ for all $g$.
A representation $\rho$ on $V\neq 0$ is \emph{irreducible} (an \emph{irrep}) if its only invariant subspaces are $0$ and $V$, and \emph{completely reducible} if $V$ is a direct sum of irreducible invariant subspaces.
\end{Definition}

Irreps are the atoms of representation theory, and the intertwiners between atoms are strongly constrained.

\begin{Lemma}[Schur]\label{lem:schur}
Let $\rho_1,\rho_2$ be irreducible and $T : V_1\to V_2$ an intertwiner. Then
(1) either $T = 0$ or $T$ is an isomorphism, so that $\operatorname{Hom}_G(V_1,V_2) = 0$ if $\rho_1\not\cong\rho_2$;
(2) over $\mathbb{C}$, every intertwiner $T : V\to V$ of an irrep with itself is a scalar, $T = \lambda\operatorname{id}_V$.
\end{Lemma}
\begin{proof}[Proof sketch]
(1) If $v\in\ker T$ then $T\rho_1(g)v = \rho_2(g)Tv = 0$, so $\ker T$ is invariant, hence $0$ or $V_1$; likewise $\operatorname{im}T$ is $0$ or $V_2$.
(2) $T$ has an eigenvalue $\lambda\in\mathbb{C}$; $T-\lambda\operatorname{id}_V$ is an intertwiner with nonzero kernel, hence zero by (1).
\end{proof}

Over $\mathbb{R}$, (2) weakens: $\operatorname{Hom}_G(V,V)$ is isomorphic to $\mathbb{R}$, $\mathbb{C}$ or the quaternions.
For example, the rotations $SO(2)$ acting on $\mathbb{R}^2$ commute with every rotation-scaling $\bigl(\begin{smallmatrix}a&-b\\b&a\end{smallmatrix}\bigr)$.
This case must be handled explicitly when building real-valued steerable networks \citep{weiler2019anips}.


% Lie groups and Lie algebras
\subsection{Lie Groups and Lie Algebras}
\label{subsubsec:lie-groups}

The groups of \Cref{ex:groups} are either \emph{discrete}, with no notion of nearby elements ($S_n$, $C_n$, $D_n$, $\mathbb{Z}^n$), or \emph{continuous}, with elements varying smoothly ($SO(n)$, $\mathbb{R}_{>0}$).
Discrete groups are handled combinatorially; continuous groups can in addition be differentiated.

\begin{Definition}[Lie group]\label{def:lie-group}
A \emph{Lie group} is a group $G$ that is also a smooth manifold, such that multiplication $(g,h)\mapsto gh$ and inversion $g\mapsto g^{-1}$ are smooth.
\end{Definition}

\begin{Example}[Lie groups]\label{ex:lie-groups}
$GL(n,\mathbb{R})$ is an open subset of $\mathbb{R}^{n^2}$ ($\det\neq 0$), hence a Lie group of dimension $n^2$; $O(n)$ and $SO(n)$ have dimension $n(n-1)/2$; $E(n)$ and $SE(n)$ have dimension $n(n+1)/2$; $\mathbb{R}_{>0}$ has dimension $1$.
The \emph{Lorentz group} $O(1,3)$ is the group of linear maps of $\mathbb{R}^4$ preserving the Minkowski form $-t^2+x^2+y^2+z^2$, used in particle-physics architectures (\Cref{sec:gdl}).
Finite groups are $0$-dimensional Lie groups.
\end{Example}

For a smooth action of a Lie group, stabilizers $G_x$ are closed Lie subgroups, and \Cref{thm:orbit-stabilizer} becomes a statement about manifolds.
If the action is proper (e.g.\ $G$ compact), $G\cdot x\cong G/G_x$ is an embedded submanifold and
\[
    \dim(G\cdot x) \;=\; \dim G - \dim G_x .
\]
For the layer symmetry $\mathbb{R}^n_{>0}\rtimes S_n$ of \Cref{ex:semidirect-products}, a configuration in which no row of $W_1$ vanishes has a finite stabilizer, since $D_\lambda W_1 = W_1$ forces $\lambda = \mathbf{1}$.
Its orbit is therefore an $n$-dimensional submanifold of parameter space on which the network function is constant.
This dimension count underlies the analysis of identifiability and functional dimension (\Cref{sec:parameter-identifiability}).

A Lie group is described to first order by its tangent space at the identity.

\begin{Definition}[Lie algebra]\label{def:lie-algebra}
The \emph{Lie algebra} of a Lie group $G$ is $\mathfrak{g} = T_eG$ with the Lie bracket $[\cdot,\cdot]$, a bilinear antisymmetric map satisfying the Jacobi identity $[A,[B,C]]+[B,[C,A]]+[C,[A,B]] = 0$.
For a matrix group $G\leq GL(n,\mathbb{R})$, $\mathfrak{g} = \{\gamma'(0) : \gamma \text{ a smooth curve in } G,\ \gamma(0) = I\}$ is a linear subspace of $\operatorname{Mat}_n(\mathbb{R})$ and $[A,B] = AB-BA$.
\end{Definition}

\begin{Example}[Matrix Lie algebras]\label{ex:matrix-lie-algebras}
Differentiating $\gamma(t)^\top\gamma(t) = I$ at $t=0$ gives $A^\top + A = 0$, so $\mathfrak{so}(n)$ is the space of antisymmetric matrices, of dimension $n(n-1)/2 = \dim SO(n)$.
For $n=2$ it is spanned by $J = \bigl(\begin{smallmatrix}0&-1\\1&0\end{smallmatrix}\bigr)$, and $\exp(\theta J)$ is the rotation by $\theta$.
Similarly $\mathfrak{gl}(n,\mathbb{R}) = \operatorname{Mat}_n(\mathbb{R})$, and the Lorentz algebra is $\{A : A^\top\eta = -\eta A\}$ with $\eta = \operatorname{diag}(-1,1,1,1)$.
\end{Example}

\begin{Definition}[Exponential map]\label{def:exp-log-map}
For a matrix group, $\exp : \mathfrak{g}\to G$, $\exp(A) = \sum_{k\geq 0}A^k/k!$, and $t\mapsto\exp(tA)$ is the unique \emph{one-parameter subgroup} ($\exp((s+t)A) = \exp(sA)\exp(tA)$) with velocity $A$ at $t=0$.
$\exp$ is a diffeomorphism from a neighbourhood of $0$ onto a neighbourhood of $e$, with local inverse $\log$.
It is surjective for compact connected $G$, but not in general (e.g.\ for $SL(2,\mathbb{R})$).
\end{Definition}

\begin{Definition}[Infinitesimal generators]\label{def:infinitesimal-generators}
Given a representation $\rho : G\to GL(V)$ of a Lie group, its \emph{infinitesimal generators} are
\[
    d\rho(A) \;=\; \frac{d}{dt}\Big|_{t=0}\rho\bigl(\exp(tA)\bigr)\in\operatorname{End}(V), \qquad A\in\mathfrak{g}.
\]
The map $d\rho : \mathfrak{g}\to\operatorname{End}(V)$ is a Lie algebra representation, $d\rho([A,B]) = [d\rho(A),d\rho(B)]$, and $\rho(\exp A) = \exp(d\rho(A))$.
\end{Definition}

Since a connected Lie group is generated by $\exp(\mathfrak{g})$, a nonlinear constraint over all of $G$ reduces to a linear one over a basis $A_1,\dots,A_k$ of $\mathfrak{g}$.

\begin{Proposition}[Linearized equivariance]\label{prop:lie-algebra-equivariance}
Let $G$ be a connected Lie group with representations $\rho, \rho'$ on $V,V'$. A linear map $T : V\to V'$ is $G$-equivariant if and only if
\[
    d\rho'(A_i)\,T - T\,d\rho(A_i) = 0, \qquad i = 1,\dots,k.
\]
\end{Proposition}

Differentiating $\rho'(\exp tA)T = T\rho(\exp tA)$ at $t=0$ gives necessity.
Conversely, the constraint gives $\exp(d\rho'(A))T = T\exp(d\rho(A))$ term by term, which is equivariance on the generating set $\exp(\mathfrak{g})$.
For disconnected groups such as $O(n)$ one adds one constraint per component.
This linear system is solved numerically to build equivariant layers for arbitrary matrix groups \citep{finzi2021p3icml}, and its residual is the penalty of approximately equivariant models (\Cref{sec:soft-equivariance}).
For a nonlinear map $\phi$, the analogous first-order quantity is the \emph{Lie derivative}
\[
    (\mathcal{L}_A\phi)(x) = \frac{d}{dt}\Big|_{t=0}\rho_Y(\exp tA)^{-1}\,\phi\bigl(\rho_X(\exp tA)\,x\bigr),
\]
which vanishes for all $A\in\mathfrak{g}$ iff $\phi$ is equivariant under connected $G$.
Learning $\mathfrak{g}$, i.e.\ generators $A_i$ or $d\rho(A_i)$, rather than the group itself is the basis of most continuous symmetry-discovery methods (\Cref{sec:sym-disc-from-data}).
One-parameter subgroups acting on parameters are also the symmetries from which conservation laws of gradient flow follow (\Cref{sec:conservation-laws}).


% Compactness and Peter-Weyl theorem guaranteeing reducibility, non-compact groups are more difficult: no reducibility of representations, Haar measure
\subsection{Compactness, Haar Measure and the Peter--Weyl Theorem}
\label{subsec:compactness-peter-weyl}

Whether an equivariant problem admits a systematic solution depends mostly on whether the group is compact.

\begin{Definition}[Compact group]\label{def:compact-group}
A Lie (more generally, topological) group is \emph{compact} if it is compact as a topological space.
A matrix group is compact iff it is closed and bounded in $\mathbb{R}^{n^2}$.
\end{Definition}

\begin{Example}[Compact and non-compact groups]\label{ex:compactness}
Finite groups, $O(n)$ and $SO(n)$ are compact; for $O(n)$, $A^\top A = I$ gives $\sum_{ij}A_{ij}^2 = n$.
$GL(n,\mathbb{R})$, $\mathbb{R}_{>0}$, the translations $\mathbb{R}^n$, $E(n)$, $SE(n)$ and the Lorentz group $O(1,3)$ are not: they contain unbounded sequences (dilations $\lambda\to\infty$, translations, boosts).
\end{Example}

\subsubsection{Haar measure and group averaging}
\label{subsubsec:haar-measure}

\begin{Theorem}[Haar measure]\label{thm:haar-measure}
Every locally compact group $G$ (in particular every Lie group) carries a nonzero Radon measure $\mu$, unique up to a positive scalar, that is \emph{left-invariant}: $\mu(gS) = \mu(S)$ for all $g\in G$ and Borel $S\subseteq G$.
\end{Theorem}

For a finite group $\mu$ is the counting measure, for $\mathbb{R}^n$ the Lebesgue measure, and for $SO(2)$ the angle measure $d\theta$.

\begin{Proposition}\label{prop:compact-haar-finite}
If $G$ is compact, then $\mu(G)<\infty$, and $\mu$ can be normalized to a probability measure that is both left- and right-invariant. If $G$ is non-compact, $\mu(G) = \infty$.
\end{Proposition}

A finite Haar measure makes it possible to average over the group.
For compact $G$ and any map $\phi:\mathcal{X}\to\mathcal{Y}$ with $\rho_Y$ linear, the \emph{group average} (or Reynolds operator)
\[
    (\mathcal{R}\phi)(x) \;=\; \int_G \rho_Y(g)^{-1}\,\phi\bigl(\rho_X(g)\,x\bigr)\,d\mu(g)
\]
is $G$-equivariant, and $\mathcal{R}\phi = \phi$ iff $\phi$ is already equivariant.
Equivariance of $\mathcal{R}\phi$ follows from the substitution $g\mapsto gh$ and right-invariance of $\mu$.
Thus $\mathcal{R}$ is a projection onto equivariant maps.
It is the continuous form of the orbit averaging behind data augmentation, feature averaging and the symmetrization of data densities (\Cref{sec:data-augmentation,sec:symmetry-breaking-dl}).
For non-compact $G$ the integral diverges, and averaging must be restricted to a compact subset or subgroup.

\subsubsection{The Peter--Weyl theorem}
\label{subsubsec:peter-weyl}

\begin{Theorem}[Peter--Weyl]\label{thm:peter-weyl}
Let $G$ be a compact group with normalized Haar measure $\mu$, and let $\widehat{G}$ denote the set of equivalence classes of irreducible unitary representations of $G$. Then:
\begin{enumerate}
    \item every irrep of $G$ is finite-dimensional, and every continuous unitary representation on a Hilbert space is a (Hilbert) direct sum of irreps;
    \item every finite-dimensional representation is equivalent to a unitary one, hence completely reducible;
    \item $L^2(G,\mu)\cong\bigoplus_{\pi\in\widehat{G}}V_\pi\otimes V_\pi^*$, i.e.\ the matrix coefficients $g\mapsto\pi(g)_{ij}$ of the irreps span a dense subspace of $L^2(G,\mu)$, and each $\pi$ occurs in the regular representation with multiplicity $d_\pi = \dim V_\pi$.
\end{enumerate}
\end{Theorem}

Part (2) follows from averaging.
For any inner product on $V$, $\langle v,w\rangle_G = \int_G\langle\rho(g)v,\rho(g)w\rangle\,d\mu(g)$ is finite by \Cref{prop:compact-haar-finite} and $G$-invariant.
The orthogonal complement of an invariant subspace is then invariant, and induction on $\dim V$ gives complete reducibility.

\subsubsection{Isotypic decomposition and equivariant layers}
\label{sec:isotypic-decomposition}

\begin{Definition}[Isotypic decomposition]\label{def:isotypic-decomposition}
Let $\rho$ be a finite-dimensional representation of a compact group on $V$.
Grouping its irreducible summands by type gives
\[
    V \;=\; \bigoplus_{\pi\in\widehat{G}} V_\pi, \qquad V_\pi\cong\pi^{\oplus m_\pi}\cong\pi\otimes\mathbb{C}^{m_\pi},
\]
where the \emph{multiplicity} $m_\pi\geq 0$ counts the copies of $\pi$, and $G$ acts trivially on the \emph{multiplicity space} $\mathbb{C}^{m_\pi}$.
The \emph{isotypic components} $V_\pi$ are canonical; their splitting into individual copies of $\pi$ is not.
\end{Definition}

\begin{Proposition}[Equivariant linear maps are block-diagonal]\label{prop:schur-blocks}
Let $V = \bigoplus_\pi\pi\otimes\mathbb{C}^{m_\pi}$ and $V' = \bigoplus_\pi\pi\otimes\mathbb{C}^{m'_\pi}$ (over $\mathbb{C}$).
Every intertwiner $T : V\to V'$ maps $V_\pi$ into $V'_\pi$ and has the form
\[
    T\big|_{V_\pi} = \operatorname{id}_\pi\otimes A_\pi, \qquad A_\pi\in\mathbb{C}^{m'_\pi\times m_\pi}.
\]
Hence $\dim\operatorname{Hom}_G(V,V') = \sum_\pi m_\pi m'_\pi$, compared with $\dim V\cdot\dim V'$ for unconstrained linear maps.
\end{Proposition}

This is \Cref{lem:schur} applied block by block: part (1) kills the blocks between inequivalent irreps, and part (2) makes each block between two copies of $\pi$ a scalar.
The resulting parameter reduction is the mechanism of steerable CNNs.
There, feature channels are organized into irreps of $O(2)$, $SO(3)$ or their finite subgroups, and the kernel constraint is solved independently for each pair of irreps \citep{weiler2018anips,weiler2019anips,cesa2021}.
For $SO(3)$, whose irreps $D_\ell$ ($\ell = 0,1,2,\dots$) have dimension $2\ell+1$, the tensor product of two irreps decomposes as
\[
    D_{\ell_1}\otimes D_{\ell_2}\;\cong\;\bigoplus_{\ell = |\ell_1-\ell_2|}^{\ell_1+\ell_2}D_\ell .
\]
The change of basis realizing this isomorphism is given by the \emph{Clebsch--Gordan coefficients}.
It is the bilinear operation by which tensor field networks and \texttt{e3nn} combine features of different degrees while remaining equivariant \citep{thomas2018,geiger2022}, and its cost motivates the optimized tensor products of \Cref{sec:gdl}.

\subsubsection{Non-compact groups}
\label{subsubsec:noncompact-representations}

For non-compact $G$, $\mu(G) = \infty$, the averaging argument fails and so can both conclusions of \Cref{thm:peter-weyl}.
First, complete reducibility fails.
$\mathbb{R}$ acts on $\mathbb{R}^2$ by $\rho(t) = \bigl(\begin{smallmatrix}1&t\\0&1\end{smallmatrix}\bigr)$; the line $\mathbb{R}\times\{0\}$ is invariant, but it has no invariant complement.
Second, irreducible unitary representations of non-compact non-abelian groups are typically infinite-dimensional; for the Lorentz group, all nontrivial ones are.
The discrete sum over $\widehat{G}$ is then replaced by a direct integral against a Plancherel measure; for $G = \mathbb{R}^n$ this is ordinary Fourier analysis.
Equivariant architectures for non-compact groups therefore rely on other devices: exploiting a compact subgroup (the rotation part of $SE(n)$), using the Lie algebra constraint of \Cref{prop:lie-algebra-equivariance} directly \citep{finzi2021p3icml}, or working with finite-dimensional, non-unitary representations \citep{bogatskiy2020p3icml}.


% Harmonic analysis
\subsection{Harmonic Analysis on Groups}
\label{subsec:harmonic-analysis}

The Peter--Weyl theorem generalizes Fourier series from the circle to any compact group; in this language, equivariant linear layers are spectral filters.

\begin{Definition}[Fourier transform on a compact group]\label{def:group-fourier-transform}
For $f\in L^2(G,\mu)$ and $\pi\in\widehat{G}$ realized as $\pi: G\to U(d_\pi)$, the \emph{Fourier coefficient} of $f$ at $\pi$ is the matrix
\[
    \widehat{f}(\pi) \;=\; \int_G f(g)\,\pi(g)\,d\mu(g)\;\in\;\mathbb{C}^{d_\pi\times d_\pi},
\]
and $f$ is recovered as $f(g) = \sum_{\pi\in\widehat{G}}d_\pi\operatorname{tr}\bigl(\widehat{f}(\pi)\,\pi(g^{-1})\bigr)$ (\Cref{thm:peter-weyl}(3)).
\end{Definition}

For abelian $G$ all irreps are $1$-dimensional characters: applying \Cref{lem:schur}(2) to $\rho(h)$, which commutes with every $\rho(g)$, shows it is a scalar.
For $G = C_n$ the characters are $\pi_k(j) = e^{2\pi i jk/n}$, and \Cref{def:group-fourier-transform} is the discrete Fourier transform (up to sign and normalization conventions).

\begin{Proposition}[Convolution theorem]\label{prop:convolution-theorem}
Define the group convolution $(f*\kappa)(g) = \int_G f(h)\,\kappa(h^{-1}g)\,d\mu(h)$. Then $\widehat{f*\kappa}(\pi) = \widehat{f}(\pi)\,\widehat{\kappa}(\pi)$ for every $\pi\in\widehat{G}$.
\end{Proposition}

The map $f\mapsto f*\kappa$ commutes with left translation $(L_uf)(g) = f(u^{-1}g)$: substituting $h\mapsto uh$ gives $(L_uf)*\kappa = L_u(f*\kappa)$.
Conversely, every bounded linear map on $L^2(G)$ commuting with left translations is of this form, with $\kappa$ possibly a distribution (e.g.\ the Dirac mass at $e$ for the identity).
This is the group convolution layer of \citet{cohen2016}, and \citet{kondor2018p3icml} show that convolution is not only sufficient but also necessary for equivariance to compact groups acting on homogeneous spaces.
On the Fourier side such a layer is block-diagonal, with one matrix $\widehat{\kappa}(\pi)$ per irrep, which is \Cref{prop:schur-blocks} for the regular representation.
Spherical CNNs use this structure with a generalized FFT on $SO(3)$ \citep{cohen2018}, and bispectral networks learn the Fourier basis of an unknown compact group from data \citep{sanborn2022}.

\begin{Definition}[Spherical harmonics]\label{def:spherical-harmonics}
The space $\mathcal{H}_\ell$ of \emph{degree-$\ell$ spherical harmonics} consists of the restrictions to $S^2$ of harmonic homogeneous polynomials of degree $\ell$ on $\mathbb{R}^3$.
It has dimension $2\ell+1$, carries the irrep $D_\ell$ of $SO(3)$, and $L^2(S^2) = \bigoplus_{\ell\geq 0}\mathcal{H}_\ell$.
\end{Definition}

In an orthonormal basis $Y_\ell = (Y_\ell^m)_{m=-\ell}^{\ell}$ of $\mathcal{H}_\ell$, the defining property is
\[
    Y_\ell(Rx) = D_\ell(R)\,Y_\ell(x), \qquad R\in SO(3),\ x\in S^2 .
\]
So $x\mapsto Y_\ell(x/\lVert x\rVert)$ turns a relative position vector into an equivariant feature of type $D_\ell$.
This is how Euclidean-equivariant networks embed geometry before combining degrees with Clebsch--Gordan tensor products (\Cref{sec:isotypic-decomposition}).
The standard reference for this material is \citet{folland2016}.


% Homogeneous Spaces cohen2019anips
\subsection{Homogeneous Spaces}
\label{subsec:homogeneous-spaces}

Data often live not on a group but on a space that a group moves around transitively: a sphere, the plane, a Grassmannian.

\begin{Definition}[Homogeneous space]\label{def:homogeneous-space}
A \emph{homogeneous space} of $G$ is a set (or manifold) $\mathcal{X}$ with a transitive $G$-action.
Fixing $x_0\in\mathcal{X}$ with stabilizer $H = G_{x_0}$, \Cref{thm:orbit-stabilizer} gives $\mathcal{X}\cong G/H$ via $gH\mapsto g\cdot x_0$ (as manifolds for proper smooth actions, \Cref{subsubsec:lie-groups}).
\end{Definition}

\begin{Example}[Homogeneous spaces]\label{ex:homogeneous-spaces}
The sphere $S^{n-1} = SO(n)/SO(n-1)$, with $SO(n-1)$ the rotations fixing a pole.
The Euclidean space $\mathbb{R}^n = E(n)/O(n)$, with $O(n)$ the stabilizer of the origin; the planar pixel grid is similarly $\mathbb{Z}^2 = (\mathbb{Z}^2\rtimes C_4)/C_4$.
The Grassmannian $\operatorname{Gr}(k,n) = O(n)/(O(k)\times O(n-k))$ of $k$-planes in $\mathbb{R}^n$.
A $G$-torsor (\Cref{def:free-transitive}) is the case $H=\{e\}$.
\end{Example}

Functions on $G/H$ are the functions on $G$ that are invariant under right translation by $H$, so harmonic analysis on $G/H$ is that of $G$ restricted to the irreps $\pi$ containing $H$-fixed vectors.
For $S^2 = SO(3)/SO(2)$ every $D_\ell$ has a one-dimensional space of $SO(2)$-fixed vectors, which recovers $L^2(S^2) = \bigoplus_\ell\mathcal{H}_\ell$.
\citet{cohen2019anips} show that equivariant linear maps between feature fields on $G/H$ are exactly convolutions with kernels $\kappa : G\to\operatorname{Hom}(V_{\mathrm{in}},V_{\mathrm{out}})$ satisfying
\[
    \kappa(h\,g\,h') \;=\; \rho_{\mathrm{out}}(h)\,\kappa(g)\,\rho_{\mathrm{in}}(h') \qquad \text{for all } h,h'\in H.
\]
Here $\rho_{\mathrm{in}},\rho_{\mathrm{out}}$ are representations of the stabilizer $H$.
For $\mathcal{X} = \mathbb{R}^n = E(n)/O(n)$ this reduces to the steerable kernel constraint
\[
    \kappa(Rx) = \rho_{\mathrm{out}}(R)\,\kappa(x)\,\rho_{\mathrm{in}}(R)^{-1}, \qquad R\in O(n),
\]
which \Cref{prop:schur-blocks} solves irrep by irrep \citep{weiler2018anips,weiler2019anips}.
Ordinary CNNs are the case $G = \mathbb{Z}^2$, $H = \{e\}$.


% Manifold/Gauge Theory: fiber bundles, principal bundle, section, steerable kernels, local symmetries, free action
\subsection{Fiber Bundles and Gauge Symmetry}
\label{subsec:manifold-gauge-theory}

On a general manifold, such as a curved mesh, there is usually no global symmetry group acting transitively.
What remains is a \emph{local} ambiguity: the choice of a reference frame at each point.
Fiber bundles formalize this.

\begin{Definition}[Fiber bundle]\label{def:fiber-bundle}
A \emph{fiber bundle} with fiber $F$ is a smooth surjection $\pi : E\to\mathcal{X}$ such that every $x\in\mathcal{X}$ has a neighbourhood $U$ with $\pi^{-1}(U)\cong U\times F$ compatibly with $\pi$ (a \emph{local trivialization}).
A \emph{section} is a smooth $s:\mathcal{X}\to E$ with $\pi\circ s = \operatorname{id}_{\mathcal{X}}$.
\end{Definition}

\begin{Definition}[Principal bundle]\label{def:principal-bundle}
A \emph{principal $H$-bundle} is a fiber bundle $\pi : P\to\mathcal{X}$ with a smooth right $H$-action on $P$ that is free and transitive on each fiber, so each fiber $\pi^{-1}(x)$ is an $H$-torsor; $H$ is the \emph{structure group}.
A local section $s : U\to P$ is a \emph{gauge}.
Any two gauges differ by a \emph{gauge transformation} $s' = s\cdot\gamma$ with $\gamma : U\to H$.
\end{Definition}

The main example is the \emph{frame bundle} of a $d$-dimensional Riemannian manifold, whose fiber over $x$ is the set of orthonormal bases of $T_x\mathcal{X}$, a principal $O(d)$-bundle.
A gauge is a local frame field, and $H = O(d)$ acts by rotating the frame independently at each point.
Local frames are unavoidable: on $S^2$ there is no continuous nonvanishing tangent vector field (hairy ball theorem), hence no global frame.

Given a representation $\rho$ of $H$ on $V$, a \emph{feature field} of type $\rho$ is a section of the \emph{associated bundle} $P\times_\rho V = (P\times V)/\!\sim$, where $(p\cdot h,v)\sim(p,\rho(h)v)$.
In a gauge $s$ it is written as a function $f: U\to V$, and a change of gauge $s' = s\cdot\gamma$ transforms it by
\[
    f'(x) \;=\; \rho\bigl(\gamma(x)\bigr)^{-1} f(x).
\]
Scalars ($\rho$ trivial) are unaffected, while tangent vectors ($\rho$ the defining representation) rotate with the frame.
A convolution on $\mathcal{X}$ is independent of the arbitrary choice of gauge iff its kernel $K : \mathbb{R}^d\to\operatorname{Hom}(V_{\mathrm{in}},V_{\mathrm{out}})$, applied in normal coordinates, satisfies
\[
    K(h^{-1}v) \;=\; \rho_{\mathrm{out}}(h)^{-1}\,K(v)\,\rho_{\mathrm{in}}(h) \qquad \text{for all } h\in H.
\]
This is the steerable constraint of \Cref{subsec:homogeneous-spaces}, now imposed pointwise with respect to the structure group rather than globally \citep{cohen2019p3icml,weiler2021,weiler2026}.
Symmetry thus acts in two ways: \emph{global} symmetry, one group acting on the whole of $\mathcal{X}$ (\Cref{subsec:group-actions,subsec:homogeneous-spaces}), and \emph{local} or \emph{gauge} symmetry, a group acting independently on the fiber over each point, with equivariance becoming a condition of consistency under local changes of frame.
% ===== END copy of main.tex, lines 335-805 =====


\section{Mathematical background 25/09}



\bibliographystyle{plainnat}
\bibliography{../bibliography}

\end{document}
